
Suppose we want to model an exchangeable sequence of random variables.
Let $S_n$ denote the sequence of the first $n$ variables: $S_n = (\seqExp{X})$.
As discussed earlier, exchangeability can be interpreted as a form of partial knowledge, which allows us to
represent our uncertainty using a central distribution over $S_n$.

However, this central distribution generally depends on the sequence length $n$.
According to the principle of objectivity, our beliefs should not depend on the specifics of our experimental setup,
and in particular, should not vary with the length of the sequence.

It is important to note that a distribution defined over a longer sequence induces marginal distributions over its
shorter subsequences.
This observation suggests that, in order to satisfy the principle of objectivity, we may instead determine a central
distribution over an infinite exchangeable sequence, from which the distributions of all finite initial segments can
be derived.

This requires constructing a central distribution that is well-defined over $S_{\infty}$.

\begin{definition}
    \label{def:infseq-central}
    Let \infSeq{S} be a sequence of random variables and \setQ be a set of probability measures,
    where each element of \setQ defines a distribution over each $S_n$.
    Assume \cDist[n] is a central distribution for $S_n$.
    Let \cDist be a probability measure that defines a distribution over every $S_n$.
    \cDist is a central distribution of \setQ for \infSeq{S}, if we have
    \begin{equation*}
        \toInf \kl[{\cDist[n]}][\cDist][S_n] = 0 \cdot
    \end{equation*}
\end{definition}

Intuitively, this definition states that as $n \to \infty$ the distribution that $C$ induces over $S_n$ becomes
increasingly similar to a central distribution for $S_n$.
In other words, \cDist is a central distribution for $S_{\infty}$.

Our next goal is to extend Theorem~\ref{th:estimator-to-var} so that it also applies in this asymptotic setting,
thereby providing a characterization of central distributions for infinite sequences.

\begin{theorem}
    \label{th:estimator-to-infseq}
    Let \infSeq{S} and \infSeq{T} be two sequences of random variables.
    Let \setQ be a set of probability measures, where each \inSet{Q} defines a joint distribution
    over $(S_n, T_n)$.
    Assume there exists a natural number $M$ and a fixed conditional distribution $R$
    such that, for all \inSet{Q} and $n \geq M$,
    \[
        Q(S_n, T_n) = R(S_n \mid T_n)Q(T_n).
    \]


    Let $C$ be a central distribution of \setQ for \infSeq{T} (Definition~\ref{def:infseq-central}).
    Define a probability measure $C^*$ by
    \[
        C^*(S_n, T_n) = R(S_n \mid T_n)C(T_n), \quad n \geq M.
    \]
    Then $C^*$ is a central distribution for \infSeq[]{(S_n, T_n)}.
\end{theorem}

\begin{proof}
    By Theorem~\ref{th:estimator-to-var}, for each $n \geq M$ we have
    \[
        C_n(S_n, T_n) = R(S_n \mid T_n) C_n(T_n),
    \]
    where $C_n(T_n)$ is a central distribution for $T_n$, and $C_n(S_n,  T_n)$ is a central distribution for
    $(S_n, T_n)$.
    Therefore,
    \begin{align*}
        \kl[{C_n}][C^*][S_n, T_n] &= \kl[{R(S_n \mid T_n)C_n}][R(S_n \mid T_n)C][T_n] \\
        &= \kl[{C_n}][C][T_n].
    \end{align*}
    Since $C$ is a central distribution for \infSeq{T} we have
    \[
        \toInf \kl[{C_n}][C][T_n] = 0.
    \]
    It follows that
    \[
        \toInf \kl[{C_n}][C^*][S_n, T_n] = 0,
    \]
    so $C^*$ is a central distribution for \infSeq[]{(S_n, T_n)}.
\end{proof}

We now present an asymptotic version of Theorem~\ref{th:lambda-centrality}, applicable to infinite sequences of
random variables.

\begin{theorem}
    \label{th:limiting-centrality}
    Let \infSeq{S} be a sequence of random variables.
    Suppose \cDist is a weighted average of elements of \setQ with strictly positive weights,
    and there exists a sequence of constants \infSeq{\lambda} such that
    \[
        \toInf  \sup_{\inSet{Q}} \left|\kl[][\cDist][S_n] - \lambda_n \right| = 0,
    \]
    then \cDist is a central distribution of \setQ for \infSeq{S}.
\end{theorem}

\begin{proof}
    For the purpose of this proof, we additionally assume that for each \(S_n\), there exists
    a central distribution $\cDist_n$ that satisfies the conditions of Theorem~\ref{th:lambda-centrality}.
    That is, there exists a constant \(\mu_n\) such that for every \inSet{Q}
    \[
        \kl[][\cDist_n][S_n] = \mu_n .
    \]

    For each \(n\), by the definition of central distribution, we have
    \begin{equation}
        \sup_{\inSet{Q}} \kl[][\cDist][S_n] \geq \mu_n,
        \label{eq:th_lim_central_proof_sup}
    \end{equation}
    and, by Lemma~\ref{lm:imperfection}, it also follows that
    \begin{equation}
        \inf_{\inSet{Q}} \kl[][\cDist][S_n] \leq \mu_n .
        \label{eq:th_lim_central_proof_inf}
    \end{equation}
    On the other hand, for each $n$ we have
    \[
        \left| \sup_{\inSet{Q}} \kl[][\cDist][S_n] - \lambda_n \right|
        \leq \sup_{\inSet{Q}}  \left|\kl[][\cDist][S_n] - \lambda_n \right|,
    \]
    and
    \[
        \left|\inf_{\inSet{Q}} \kl[][\cDist][S_n] - \lambda_n \right|
        \leq \sup_{\inSet{Q}}  \left|\kl[][\cDist][S_n] - \lambda_n \right|.
    \]
    This implies that the left hand sides of
    Equations~\ref{eq:th_lim_central_proof_sup} and~\ref{eq:th_lim_central_proof_inf}  both converge to $\lambda_n$.
    Now, by the squeeze theorem, the sequence \infSeq[]{\mu_n - \lambda_n} must converge to $0$.
    Therefore, for any arbitrary \(Q \in \setQ\), we have
    \begin{align*}
        \toInf \sum_{s_n} Q(s_n) \ln \frac{\cDist_n(s_n)}{\cDist(s_n)} &=
        \toInf (\kl[][\cDist][S_n] - \kl[][\cDist_n][S_n]) \\
        &= \toInf (\kl[][\cDist][S_n] -\mu_n + \lambda_n - \lambda_n) \\
        &= \toInf(\kl[][\cDist][S_n] - \lambda_n) - \toInf(\mu_n - \lambda_n) \\
        &= 0 .
    \end{align*}
    Since $\cDist_n$ satisfies the conditions of Theorem~\ref{th:lambda-centrality},
    $\kl[\cDist_n][\cDist][S_n]$ is a weighted average of
    $\sum_{s_n} Q(s_n) \ln \frac{\cDist_n(s_n)}{\cDist(s_n)} $, and we can conclude
    \[
        \toInf \kl[\cDist_n][\cDist][S_n] = 0 .
    \]
    Therefore, \cDist is a central distribution of \setQ for the sequence \infSeq{S}.
\end{proof}

\newcommand{\estLike}{Q \left( \est_{n} \mid \paramVal \right) }

Assume \setQ is a set of distributions that can be indexed by a parameter vector \param, and \infSeq{\est} is a
consistent estimator sequence for \param.
Let $P(\paramVal \mid\est_n = t)$ denote the posterior distribution obtained from an arbitrary smooth positive
prior.

Under regular parametric conditions, posterior distributions based on consistent estimators become asymptotically
insensitive to the choice of smooth positive prior.

Motivated by this phenomenon, we heuristically derive the asymptotic form of a central prior by considering the
posterior density evaluated at the point $t=\paramVal$.
A similar asymptotic construction also appears in~\cite{BERNARDO200517}.

Specifically, we define

\begin{equation}
    \label{eq:central-limit}
    \cPrior(\paramVal) \propto \toInf \frac{P(\paramVal \mid\est_n = t) \vert_{t = \paramVal}}{\mu_n},
\end{equation}
where \infSeq{\mu} is a sequence of positive constants chosen so that the limit exists for every value of \paramVal.

Let $k$ denote the normalizing constant converting proportionality into equality.
Since the logarithm is a continuous function, Equation~\ref{eq:central-limit} implies
\begin{equation*}
    \toInf \left(\ln \frac{P(\paramVal \mid\est_n = t) \vert_{t = \paramVal}}{\cPrior(\paramVal)}
               - \ln \mu_n + \ln k \right) = 0,
\end{equation*}

Since $\est_n$ is a consistent estimator, the sampling distribution $\estLike$ becomes increasingly concentrated
around $\est_n = \paramVal$ as $n \to \infty$.
Furthermore, under the usual regularity conditions, the posterior distribution becomes asymptotically
insensitive to the choice of smooth positive prior.
Motivated by this asymptotic prior-insensitivity, we heuristically replace the original prior with $c(\paramVal)$,
we obtain
\begin{equation*}
    \toInf \left( \sum_{t_n} Q(t_n \mid \paramVal) \ln
               \frac{Q(t_n \mid \paramVal)}{C(t_n)}
               - \ln \mu_n + \ln k \right) = 0 \cdot
\end{equation*}
Letting
\[
    \lambda_n = \ln \mu_n - \ln k,
\]
this suggests that, for each fixed \paramVal:
\begin{equation*}
    \toInf  \left( \Dkl{\estLike}{C(\est_n)} - \lambda_n \right) = 0 \cdot
\end{equation*}
Under additional regularity assumptions ensuring that this convergence is uniform over \paramVal, the assumptions of
Theorem~\ref{th:limiting-centrality} are satisfied.
Thus, the above asymptotic argument suggests that $\cPrior(\paramVal)$ acts as a central prior of \setQ for the
estimator sequence \infSeq{\est}.

A central distribution over an estimator cannot, by itself, be used to perform inference at the level of observed
outcomes.
To do so, we must convert it into a distribution over the primary variables of interest.
Theorems~\ref{th:estimator-to-var} and~\ref{th:estimator-to-infseq} provide a simple
method for this conversion.

\begin{example}[Binomial Experiment]
    \newcommand{\pv}{\paramVal}
    Consider a binomial experiment with $n$ Bernoulli trials.
    Denote the sequence of outcomes by $S_n$, and let $k$ be the number of successes observed in $S_n$

    Let \setQ be the set of distributions corresponding to different values of the success
    probability $\pv \in [0,1]$.
    We can consider $\pv$ as an index parameter for \setQ, and
    we write $Q(S_n \mid \pv)$ to denote the distribution $Q_{\pv}$ applied to $S_n$.
    Explicitly,
    \[
        Q(S_n \mid \pv) = \pv^k (1-\pv)^{n-k} \cdot
    \]

    Define $T_n$ to be the proportion of successes in $n$ trials.
    Note that $T_n$ is a consistent estimator for $\pv$.
    Therefore, we can use Equation~\ref{eq:central-limit} to find a central distribution for the sequence
    \infSeq{T}.

    Assuming a uniform prior on $\pv$, Bayes’ rule gives the posterior distribution
    \[
        P\left( \pv \, \middle| \, T_n = \tfrac{k}{n} \right) = (n+1)\binom{n}{k} \pv^k (1 - \pv)^{n - k},
    \]
    which is a $\mathrm{Beta}(k+1, n-k+1)$ distribution.
    Using Stirling’s formula, we derive its asymptotic form:
    \begin{equation*}
        P\left( \pv \, \middle| \, T_n = \tfrac{k}{n} \right) \sim
        \frac{(n+1)n^{(n+\hf)}}{\sqrt{2 \pi} (n-k)^{(n-k+\hf)} k^{(k+\hf)}} \, \pv^k (1 - \pv)^{n - k} \cdot
    \end{equation*}
    We now let $T_n \to \pv$:
    \begin{equation*}
        P\left(\pv \mid T_n = \pv \right) \sim \frac{n+1}{\sqrt{2 \pi n}} \, \pv^{-\hf}(1 - \pv)^{-\hf} \cdot
    \end{equation*}
    Therefore, if we define
    \[
        \mu_n = \frac{n+1}{\sqrt{2 \pi n}},
    \]
    the following limit exists:
    \[
        \toInf \frac{P\left(\pv \mid T_n = \pv \right)}{\mu_n} = \pv^{-\hf}(1 - \pv)^{-\hf} \cdot
    \]
    This implies that $c(\pv) \propto \pv^{-\hf}(1 - \pv)^{-\hf}$ is a central prior for \infSeq{T}.

    Now we apply Theorem~\ref{th:estimator-to-infseq} to obtain a central distribution for $\infSeq{S}$.
    Since $T_n$ is a deterministic function of $S_n$, the joint distribution $Q(S_n, T_n)$ coincides with the
    distribution of $S_n$: it assigns probability $Q(S_n)$ when $T_n$ is compatible with $S_n$, and zero otherwise.
    In other words, knowledge of $(S_n, T_n)$ is equivalent to knowledge of $S_n$.

    In the binomial experiment, all sequences with the same number of successes and failures have equal probability.
    Therefore,
    \[
        Q(S_n \mid \pv)
        = Q\!\left(S_n,\, T_n = \tfrac{k}{n} \mid \pv\right)
        = \frac{1}{\binom{n}{k}} Q(T_n \mid \pv) \cdot
    \]
    Let $C^*$ denote the central distribution of $\infSeq{S}$.
    By Theorem~\ref{th:estimator-to-infseq}, the distribution induced by $C^*$ on each $S_n$ is given by:
    \begin{equation*}
        C^*(S_n)
        %= C^*\!\left(S_n,\, T_n = \tfrac{k}{n}\right)
        = \frac{\gm{k+\hf} \gm{n-k+\hf}}{\gm{n+1}\gm{\hf}^2} \cdot
    \end{equation*}
\end{example}

When dealing with layered ignorance and attempting to determine the central distribution of a hierarchy of
distributions, as defined in Definition~\ref{def:hierarchy-of-distributions}, we typically encounter the case where the number of
elements at a given level of the hierarchy is finite.
In that setting, and under appropriate regularity conditions,
the central distribution reduces to an equally weighted average of the corresponding representative central
distributions.